\subsection{Uniform cocycle bounds on the required deep towers}
\label{subsec:rh-four-uniform-tower-bounds}

We supply the exact \texttt{Fr.Bounds} and large-degree hypotheses
of the rank-four component-support theorem. No bound in terms of
the auxiliary field degree is needed for the constant $C$.

\begin{theorem}[Uniform bounds on the original cyclotomic layers]
\label{thm:rh-four-x2}
Let $D$ be the original good datum, $\ell$ odd, and let $T_1$
be a fitted tower with $F_1=T_1.F$. There exists $C$ such that,
for every original \texttt{Le} tower $T$ satisfying
\texttt{IsLayer}$(\ell,m,F_1,T.F)$, and every field $k$ equipped
with a map from the original coefficient DVR killing its maximal
ideal, the frame $(T.\mathrm{fr}(\varnothing))$ satisfies
\texttt{Bounds}$(k,[T.F:\mathbb Q],C)$.
For every $n$ there is such a tower which is deep and has degree
at least $n$. In particular after fixing the starting tower,
one can use $n=30(C+1)^2$ in the component-support theorem.
\end{theorem}
\begin{proof}
Write $\omega=\bar\epsilon$ and denote the original four residual
diagonal characters by $\bar\chi_i$, over their original finite
residue field $k_0$. Each has finite order prime to $\ell$.
Let $F_m=F_1\mathbb Q_m$, where $\mathbb Q_m$ is the degree
$\ell^m$ subfield of the cyclotomic $\mathbb Z_\ell$-extension
of $\mathbb Q$. The original \texttt{IsLayer} condition is
precisely
\[
 G_{F_m}=G_{F_1}\cap
   \{g:\epsilon(g)^{\ell-1}\equiv1\bmod\ell^{m+1}\}.
\]
Use the original residual-kernel field
\[
 E_m=F_m\,\mathbb Q(\mu_\ell)\,
         \overline{\mathbb Q}^{\,\bigcap_i\ker\bar\chi_i}.
\]
Thus $E_m/E_0$ is a layer of the cyclotomic
$\mathbb Z_\ell$-extension of the fixed abelian number field
$E_0$. The extension $E_m/F_m$ is abelian of degree prime
to $\ell$, because it is cut out over $F_m$ by characters into
$k_0^\times$. In particular $E_m$ is CM (it contains
$\mu_\ell$, and is abelian over $\mathbb Q$).

Ferrero--Washington, \emph{The Iwasawa invariant $\mu_p$ vanishes
for abelian number fields}, Ann.\ Math.\ 109 (1979), 377--395,
gives $\mu(E_\infty/E_0)=0$. By the Iwasawa structure theorem
the unramified abelian pro-$\ell$ class-group module over the
cyclotomic extension is finitely generated over $\mathbb Z_\ell$.
It follows that the $\ell$-ranks of the ideal class groups
of the $E_m$ are uniformly bounded, say by $R$.
One direct way to see the finite-layer implication is to go
high enough in the cyclotomic extension that all primes over
$\ell$ are totally ramified in the remaining extension.
The $\ell$-Hilbert class field of any such layer is then
disjoint from the remaining cyclotomic extension over that
layer. Its compositum with the infinite layer is an unramified
abelian pro-$\ell$ extension there. Its Galois group, the
$\ell$-part of the class group of that layer, is a quotient
of the class-group module. The finitely many earlier layers
can be included by increasing $R$.

The number of places of $E_m$ over the fixed finite rational
set $S=D.S$ is bounded by a constant $g$, independent of $m$.
For $\ell$ this follows from total ramification after a fixed
finite layer. At $p\ne\ell$ the cyclotomic Frobenius is the
principal-unit part of $p$ (or its inverse in the geometric
convention); it has infinite order, since $p^{\ell-1}\ne1$
as a rational integer. It generates an open subgroup in the
rank-one group $\mathbb Z_\ell$, so the decomposition group
has finite index in the cyclotomic extension. Summing over
the primes of $E_0$ above $S$ proves the bound. It also bounds
the corresponding number of places of $F_m$, in particular
the number of ordinary places.

For a root weight $w\in\mathbb Z^2$ put
$\lambda=\overline{\Phi.\mathrm{rt}(w)}$ on $G_{F_m}$, with
$\Phi=T.\mathrm{fr}(\varnothing)$. We prove the bounds on the
literal spaces of locally constant cocycles, not just on
their cohomology. First consider a strict cocycle
$z\in\Phi.\mathrm{Str}(w)$, i.e. a normalized cocycle
killing $\Phi.N$ and a coboundary on each ordinary
decomposition group. Its restriction to $G_{E_m}$ is a
homomorphism, because $\lambda=1$ there. It kills inertia
away from $\ell$: those inertia elements have trivial
original residual diagonal and belong to $\Phi.N$.
It kills inertia at $\ell$ because the local cocycle is a
coboundary and the restricted character is trivial.
For the other choices of places use the cocycle conjugation
formula, which gives the same vanishing since the restriction
to $G_{E_m}$ is a homomorphism and $G_{E_m}$ is normal.
Consequently its restriction factors through the unramified
abelian $\ell$-class group of $E_m$. The kernel of restriction
has dimension at most one, by averaging cocycles on the
finite quotient $\Delta_m=\operatorname{Gal}(E_m/F_m)$,
whose order is prime to $\ell$. Hence
\[
                  \dim_{k_0}\Phi.\mathrm{Str}(w)\leq R+1 .
\]
This even treats the trivial root character, for which the
coboundary space is zero.

Every cocycle in $\Phi.\mathrm{Zgl}(w)$ is unramified
outside $S$: outside $S$ the residual diagonal is trivial
on inertia, so the killing condition applies. Inflation--restriction
and averaging on $\Delta_m$ identify its cohomology class
with a subspace of the $\omega\lambda^{-1}$ isotypic part of
\[
 H^1(G_{E_m,S},k_0(1))
  \simeq H^1(G_{E_m,S},\mu_\ell)\otimes_{\mathbb F_\ell}k_0.
\]
Here $S$ in the subscript of a Galois group means all places
above the original rational set. The identity of isotypic
characters follows because the restricted coefficient module
is trivial and the two twists differ by
$\lambda\omega^{-1}$. The Kummer sequence for $S$-integers gives
an exact sequence of $\Delta_m$-modules
\[
 0\longrightarrow
 \mathcal O_{E_m,S}^{\times}/(\mathcal O_{E_m,S}^{\times})^\ell
 \longrightarrow H^1(G_{E_m,S},\mu_\ell)
 \longrightarrow \operatorname{Cl}(\mathcal O_{E_m,S})[\ell]
 \longrightarrow0.
\]
The last term has dimension at most $R$. The free part of
the ordinary units has rational character the permutation
representation on the complex places of $E_m$ minus the
trivial character, by Dirichlet's unit theorem and the
equivariant logarithmic embedding. Its reduction modulo
$\ell$ has the same multiplicities for $\Delta_m$: its order
is prime to $\ell$, so the character idempotents over an
unramified coefficient DVR containing the roots of unity
of that order split every integral lattice. The dimensions
of the reductions are the ranks of those summands.
The $\ell$-primary torsion of the units contributes at most
one additional dimension. Enlarging to $S$-units adds at
most $g$ dimensions, by the valuation sequence: the quotient
by ordinary units is a subgroup of the free group on the
$S$-places, hence free of rank at most $g$.

Complex conjugation acts as the central element $c$ of
$\Delta_m$ since $F_m$ is real. Every complex place has
stabilizer $\{1,c\}$ in $\Delta_m$ and there are precisely
$d=[F_m:\mathbb Q]$ orbits. Thus the odd characters have
multiplicity zero in the free ordinary units, and each
even character has multiplicity at most $d$. For the original
even weights $0,w_{\rm Ps},-w_{\rm Ps}$, $\lambda$ is even
since the original residual $\Psi$ is even, and
$\omega\lambda^{-1}$ is odd. For every other root the
same argument gives the general upper bound $d$ for its
unit contribution. Restoring at most one dimension of
coboundaries, we obtain
\[
 \begin{aligned}
 \dim_{k_0}\Phi.\mathrm{Zgl}(w)&\le R+2+g
     &&(w=0,w_{\rm Ps},-w_{\rm Ps}),\\
 \dim_{k_0}\Phi.\mathrm{Zgl}(w)&\le d+R+2+g
     &&(\text{all }w).
 \end{aligned}
\]
In these inequalities one may retain the extra zero-at-mark
condition of the original $\mathrm{Zgl}$ throughout; dropping
it only enlarges the space.

For a selected ordinary place $v$ the local Euler characteristic
formula and local Tate duality give
$\dim H^1(G_{F_{m,v}},k_0(\lambda))
=[F_{m,v}:\mathbb Q_\ell]+\dim H^0(k_0(\lambda))
+\dim H^0(k_0(\omega\lambda^{-1}))$.
The coboundary dimension is $1-\dim H^0(k_0(\lambda))$.
Thus the local cocycle dimension is at most
$[F_{m,v}:\mathbb Q_\ell]+2$. Summing over the original
representatives of all places over $\ell$ gives at most
$d+2g$. Taking $C=R+2+2g$ proves all five fields of
\texttt{Fr.Bounds} over $k_0$.

For any larger field $k$, each locally constant cocycle has
finite image, by compactness, and its values span a finite
$k_0$-subspace of $k$. For a strict cocycle include the
finitely many local coboundary witnesses in that span.
Expansion in a $k_0$-basis writes it as a finite sum of
$k_0$-valued cocycles with coefficients in $k$, each satisfying
the same linear equations, including the mark and killing
conditions. The identical argument for the finite product
of local cocycle spaces proves scalar extension of every
bound, with no assumption that $k$ itself is finite.

Finally we construct the deep layers. Since $F_1$ is finite
and $\mathbb Q_\infty/\mathbb Q$ infinite,
$[F_m:\mathbb Q]$ tends to infinity. By \texttt{Ok.kc}
the $\ell$-primary part of $\psi$ is trivial on
$G_{F_1}\cap\ker\epsilon$. Its restriction factors continuously
through the closed subgroup $\epsilon(G_{F_1})$ of
$\mathbb Z_\ell^\times$. The $\ell$-primary character is
trivial on the finite prime-to-$\ell$ torsion of that subgroup
and factors through a finite quotient of its principal-unit
part. Hence it is trivial on $G_{F_m}$ for all sufficiently
large $m$. This is exactly \texttt{Deep}. The fields $F_m$
are abelian and totally real. The fixed $K=T_1.K$, original
ordering, representatives of the ordinary places above
$\ell$, and the mark $T_1.g_0^{\ell^m}$ give the original
\texttt{Le} tower (the relative degree divides $\ell^m$).
The \texttt{Ok} and \texttt{Fit} properties are preserved by
the odd-prime-power extension as in the preceding residual-chain
argument: witnesses may be raised to $\ell$-powers,
which preserve every inequality of prime-to-$\ell$ residual
values, and the quadratic-character exclusion extends by
the corresponding residual ratio. The relative unramified
quadratic condition persists by base change.
Taking $m$ sufficiently large gives the required degree
and all the bounds on this same layer.
\end{proof}

This proves the arithmetic supplier \texttt{x2} and the required
\texttt{layerEx} construction. The constant is fixed after the
auxiliary primes of the residual chain have been chosen, but
before the large cyclotomic layer is chosen. It is not a uniform
estimate for the final moving Selmer product, and is not used
as such.
